\documentclass{article}
\usepackage{amsmath}
\usepackage{amssymb}
\usepackage{listings}
\usepackage{color}
\usepackage{hyperref}

\definecolor{codegreen}{rgb}{0,0.6,0}
\definecolor{codegray}{rgb}{0.5,0.5,0.5}
\definecolor{codepurple}{rgb}{0.58,0,0.82}
\definecolor{backcolour}{rgb}{0.95,0.95,0.92}

\lstdefinestyle{mystyle}{
    backgroundcolor=\color{backcolour},   
    commentstyle=\color{codegreen},
    keywordstyle=\color{magenta},
    numberstyle=\tiny\color{codegray},
    stringstyle=\color{codepurple},
    basicstyle=\footnotesize\ttfamily,
    breakatwhitespace=false,         
    breaklines=true,                 
    captionpos=b,                    
    keepspaces=true,                 
    numbers=left,                    
    numbersep=5pt,                  
    showspaces=false,                
    showstringspaces=false,
    showtabs=false,                  
    tabsize=2
}
\lstset{style=mystyle}

\title{Full Trace Synthesis with Guided Constant Seeding}
\author{Experimental Report}
\date{\today}

\begin{document}

\maketitle

\begin{abstract}
This report details the experimental results of applying "Full Trace Synthesis" to a suite of multi-phase invariant generation benchmarks from the \texttt{s\_split} family. We demonstrate that while generic SyGuS learners struggle with extrapolation due to the "Gap of Constants," seeding the grammar with derived constants (from scalar evolution or affine analysis) and restricting boolean search using Model Based Projection (MBP) guards enables instantaneous synthesis of complex, multi-phase bitvector invariants.
\end{abstract}

\section{Introduction}
Counter-example Guided Inductive Synthesis (CEGIS) typically relies on short error traces. However, for multi-phase loops where behavior changes deep in the execution (e.g., after 5000 iterations), short traces fail to capture the necessary inductive invariants. Our "Full Trace" approach generates execution traces covering all phase transitions. We found that the primary bottleneck is not the trace length, but the inability of solvers to synthesize large, implicit constants required for the solution.

\section{Methodology}
\subsection{Full Trace Generation}
We execute the loop concretely (in Python or C++) to generate a trace $T = \{ (n, \mathbf{v}_n) \mid 0 \le n \le k \}$ that covers all known phase transitions.

\subsection{Seeded SyGuS Grammars}
Standard SyGuS grammars provided with a "bag of literals" found in the code often fail because the invariant requires constants that are linear combinations of the program constants (e.g., offsets $b$ in $y = ax + b$). We augment the grammar with:
\begin{itemize}
    \item \textbf{Implicit Intersection Constants}: Derived from the intersection of phase trajectories.
    \item \textbf{Affine Offsets}: Derived from accumulation variables (e.g., $12000$ from $\sum_{i=0}^{3000} 4$).
\end{itemize}

\subsection{MBP-Guided Boolean Search}
To reduce the search space for the boolean structure (guards), we restrict the boolean grammar to predicates identified by Model Based Projection (MBP) or abstract interpretation (e.g., $n < 4000$).

\section{Case Studies}

\subsection{Benchmark: \texttt{s\_split\_01}}
\paragraph{Problem Description}
A single variable loop with a phase transition at $n=5000$.
\begin{itemize}
    \item Phase 1 ($n < 5000$): $y$ increases.
    \item Phase 2 ($n \ge 5000$): $y$ remains constant.
\end{itemize}

\paragraph{Synthesis Result}
\textbf{Status: Success} (Generic Grammar).
The constant $5000$ (\texttt{\#x1388}) was present in the source code. The solver successfully synthesized:
\begin{lstlisting}[language=Lisp]
(define-fun fy ((n (_ BitVec 16))) (_ BitVec 16) 
  (ite (bvult n #x1388) #x1388 n))
\end{lstlisting}
Note: The solver found an equivalent form that satisfied the constraints, often picking specifically the constant value for the second phase.

\subsection{Benchmark: \texttt{s\_split\_48}}
\paragraph{Problem Description}
A complex loop with 4 phases and 10,000 steps.
\begin{enumerate}
    \item $0 \le n < 4000$: $y(n) = n$
    \item $4000 \le n < 5000$: $y(n) = 4n - 12000$
    \item $5000 \le n < 6000$: $y(n) = 28000 - 4n$
    \item $n \ge 6000$: $y(n) = 10000 - n$
\end{enumerate}

\paragraph{Challenge}
The generic solver timed out. The required constants $12000$ and $28000$ do not appear in the code; they are the result of integrating the differential equations of previous phases.

\paragraph{Solution}
We seeded the SyGuS grammar with the derived constants: \texttt{\#x2EE0} (12000) and \texttt{\#x6D60} (28000). We also provided the phase boundaries (4000, 5000, 6000) as grammar terminals.

\paragraph{Result}
\textbf{Status: Success} (Seeded, $<1$s).
\begin{lstlisting}[language=Lisp]
(ite (bvult n #x0fa0) n 
  (ite (and (bvult n #x1388) (bvuge n #x0fa0)) 
       (bvsub (bvmul n 4) #x2ee0) ; 4n - 12000
       ... ))
\end{lstlisting}

\subsection{Benchmark: \texttt{s\_split\_42}}
\paragraph{Problem Description}
A chained dependency where $z$ depends on $y$, and $y$ depends on $n$.
\begin{itemize}
    \item $y$ transitions at $n=1765$.
    \item $z$ transitions when $y \ge 5765$.
\end{itemize}

\paragraph{Challenge}
The condition $y \ge 5765$ maps to $n \ge 3765$ (since $y = 2n - 1765$ in that region). The constant $3765$ is completely absent from the source code and must be derived by solving $2n - 1765 = 5765$.

\paragraph{Solution}
We calculated the intersection point $3765$ (\texttt{\#x0EB5}) and seeded it into the grammar.

\paragraph{Result}
\textbf{Status: Success} (Seeded, $<1$s).
The solver correctly synthesized $z(n)$ using the hidden constant:
\begin{lstlisting}[language=Lisp]
(define-fun fz ((n (_ BitVec 16))) (_ BitVec 16) 
  (ite (bvult n #x0EB5) 
       (bvadd n n)
       (bvsub (bvmul n 3) #x0EB5))) ; 3n - 3765
\end{lstlisting}

\subsection{Benchmark: \texttt{s\_split\_05}}
\paragraph{Problem Description}
Involves a geometric progression (powers of 2) with a phase transition.
\begin{itemize}
    \item Phase 1: Constant value.
    \item Phase 2: Doubles every step ($z_{i+1} = z_i \cdot 2$).
\end{itemize}

\paragraph{Challenge}
This benchmark highlighted two critical limitations:
\begin{enumerate}
    \item \textbf{Trace Length}: The initial trace generation (4 steps) was insufficient. The solver could not distinguish between linear and exponential growth, producing incorrect candidates.
    \item \textbf{Grammar}: Standard BV grammars struggle to synthesize $2^n$ via repeated multiplication.
\end{enumerate}

\paragraph{Solution}
We employed a three-pronged approach:
\begin{enumerate}
    \item \textbf{MBP Seeding}: Extracted the phase transition guard automatically using Model Based Projection.
    \item \textbf{Trace Extrapolation}: Extended the trace length to 6 steps to capture the divergence of the geometric series.
    \item \textbf{Grammar Augmentation}: Seeded the grammar with the shift operator \texttt{(bvshl 1 Start)}.
\end{enumerate}

\paragraph{Result}
\textbf{Status: Success} (Enhanced Trace + Grammar).
The solver correctly synthesized the exponential function using bit-shifts:
\begin{lstlisting}[language=Lisp]
(define-fun fz ((n (_ BitVec 16))) (_ BitVec 16) 
  (ite (bvult n #x0001) 
       #x0001 
       (bvshl #x0001 (bvsub n #x0001)))) ; 2^(n-1)
\end{lstlisting}

\section{Discussion}
\subsection{The Necessity of Constant Seeding}
Our experiments confirm that "Neural" or "Enumerative" synthesis is computationally strictly bound by the "Manhattan Distance" to the target constants in the operator graph. If a constant like $28000$ requires 5+ additions/multiplications to construct from inputs, the search depth becomes prohibitive. Pre-calculating these constants via lightweight static analysis (e.g., scalar evolution) is not just an optimization, but a prerequisite for success.

\subsection{Bootstrap Invariant Seeding}
We implemented a \texttt{--sygus-invariants} flag that extracts candidate invariants from CHC bodies via a lightweight seed-mining approach. For bitvector CHCs, this requires:
\begin{enumerate}
    \item \textbf{BV to LIA Translation}: Convert bitvector constraints to linear integer arithmetic using \texttt{Bv2LiaTranslator}.
    \item \textbf{Invariant Extraction}: Collect comparison constraints from CHC bodies that only mention invariant variables.
    \item \textbf{LIA to BV Translation}: Convert extracted invariants back to bitvector expressions using \texttt{Lia2BvTranslator}.
\end{enumerate}

For \texttt{s\_split\_05}, this automatically discovered:
\begin{itemize}
    \item \texttt{bvult(\_FH\_1, 0)} (phase condition: $y < 0$)
    \item \texttt{\_FH\_2 = 1} (initial value of $z$)
    \item \texttt{bvugt(\_FH\_0, 0)} ($x > 0$ invariant)
\end{itemize}

These invariants are added to both the boolean grammar (as guard conditions) and the constant pool (extracting numeric literals from the invariants).

\subsection{BitVector Overflow Safety}
In \texttt{s\_split\_42}, we observed that synthesizers trained on integer abstractions fail validation on 16-bit BitVector benchmarks due to signed overflow. Synthesized invariants must be validated against bit-precise traces that faithfully model wrap-around behavior.

\section{Conclusion}
Full Trace Synthesis is a viable strategy for multi-phase loop invariant generation, provided it is coupled with:
\begin{enumerate}
    \item \textbf{Constant Analysis}: To discover implicit constants and phase intersections.
    \item \textbf{MBP/Guard Seeding}: To guide the boolean structure search.
\end{enumerate}
When these conditions are met, CVC5 can synthesize complex non-linear and multi-phase invariants in under 1 second.

\end{document}
